\documentclass[conference]{IEEEtran}
\IEEEoverridecommandlockouts
% The preceding line is only needed to identify funding in the first footnote. If that is unneeded, please comment it out.

% Packages
\usepackage{cite}
\usepackage{amsmath,amssymb,amsfonts}
\usepackage{algorithmic}
\usepackage{graphicx}
\usepackage{textcomp}
\usepackage{xcolor}

\def\BibTeX{{\rm B\kern-.05em{\sc i\kern-.025em b}\kern-.08em
    T\kern-.1667em\lower.7ex\hbox{E}\kern-.125emX}}


\begin{document}

\title{Finite Element approximation of the Elliptic Eigenproblems: Convergence Analysis and Numerical Experiments}
\author{}

\maketitle
\section{Introduction}
In modern science to solve complex problems various simulations based on mathematical models are used. By doing so, researchers are able to quickly test hypotheses, change parameters and solve equations that are too complex or impossible to physically recreate. One of these mathematical tools is Finite Element Method (FEM) that is used to solve differntial equations.

Differential equations are widely used for predictions models and simulations of the physical phenomenas. However, the closer they describe the reality, the more complex they become, therefore scientists reduce the number variables to the be left with the most useful ones. But, even reduced system can be too complex to solve it analytically, that is the reason why numerical analysis comes in. it allows to design and implememt an algorithms to find approximate but highly accurate numeric solutions. 

So, to solve the elliptic partial differential problems we will use the numerical method which finds the solutions accurately, however still has the error, so in this project the finite element approximation will be explored and will be checked how accurately it gives the solution and how quickly the solution goes to the desirable value (convergence analysis) by comparing with the known previously obtained results.

\section{Mathematical Background}

\subsection{Elliptic partial differential equations}
One of the branches of differential equations is elliptic differential equation, which has the general form:
\begin{equation}
A*u_{xx} + B*u_{xy} + C*u_{yy} + D*u_{x} + E*u_{y} + F*u = G (1)
\end{equation}

The equation (1) is classified as elliptic if the discriminant is strictly negative:

\begin{equation}
B^2 - 4*A*C < 0
\end{equation}

Elliptic equations differ from hyperbolic and parabolic equations because they describe the steady state, time-independent phenomena. Moreover, their solutions depend only on spatial variables and require boundary conditions rather than initial conditions, the most common are Dirichlet (specifying the value on the boundary) and Neumann (specifying the derivative or the flux at the boundary). In addition to this, changes at any point of the domain influences the solution, reflecting the global nature of elliptic problems.


\subsection{Eigenvalue problems}

\subsection{Weak formulation}
- weak formulation also known as Galerkin method, is an approach of multiplying the problem by an arbitrary test function that will disappear at the end
- using integration by parts (for 1D) or Green's formula (for 2D) will obtain the weak form


\subsection{Bilinear forms}

\subsection{Galerkin approximation}

\subsection{Finite-dimensional spaces}
-instead of dealing with the u(x, y) we can use a discrete function that we want to find
-Finite element discretisation is the process of breaking down a complex continuous physical domain into a finite number of smaller, interconnected subregions (elements). This transforms complex governing differential equations into a manageable system of algebraic equations that can be solved numerically
-Elements & Nodes: The domain is divided into a mesh of shapes like triangles and other. The vertices and designated points on these shapes are called nodes

Local Approximation: While the exact physical behavior of a structure is continuous, the solution within each individual element is approximated using simple mathematical functions, typically linear or quadratic polynomials.
Continuity: Shape functions, ensure that the solution remains continuous across the boundaries between adjacent elements.


\subsection{Generalised matrix eigenvalue problem}

\section{Finite Element Method (FEM)}

Finite element method a general tool for numerical solution of two-point boundary value problems
- it can easily deal with complex geometries and higher order approximations of the solutions
- basic idea is to divide the domain into elements (in our case triangles) and look for a polynomial approximations to the unknown functions on each triangle
- rewrite the boundary value problem as a variational equation and then seek a solution approximation of this equation from the space of continuous piecewise polynomials

problem 

\subsection{Mesh}


\subsection{Basis functions}

\subsection{Finite element approximation}
- the aim is to approximate the unknown function by a sum 
$u(x) = \sum_{i=0}^N c_i*\psi_i(x)$,
where $\psi_i$ are prescribed functions and $c_0,\dots, c_N$ are unknown coefficients to be determined.


\subsection{Piecewise polynomial spaces}


\subsubsection{P1 Elements}
- p1 elements are the continuous piecewise polynomials of degree 1, therefore their gradients are the constants
- have the theoretical $O(h^2)$


\subsubsection{P2 Elements}
- p2 elements are the continuous piecewise polynomials of degree 2, therefore their gradients will be linear
- give better approximation $O(h^4)$ in other words will give us more precise approximation with less mesh refinement

\subsection{Finite element approximation}
- the aim is to approximate the unknown function by a sum 
$u(x) = \sum_{i=0}^N c_i*\psi_i(x)$,
where $\psi_i$ are prescribed functions and $c_0,\dots, c_N$ are unknown coefficients to be determined.


\section{Elliptic Eigenproblems}

\subsection{Laplace Eigenproblem}

\subsection{
    \text{Schr\"{o}dinger's Eigenproblem}
}

\section{Goals and Objectives}
- research about the FEM
- implement fem on practice for 1D and then 2D
- analyse the results
- interpret the results
- write a report


\section{Implementation}
- to complete this project, the VS code was used to write and run the code in python
- for advanced visualisations Paraview was also used


\subsection{Mesh generation}
- the mesh was generated on the unit square and in the final version a person will be able to choose the number of points per edge of the unit square
- after generating the square with nodes the square, which is our domain, was divided into the triangulations using the Delaunay method.
- to avoid breaking of the code all the nodes were sorted in ascending order


\subsection{Local element matrices}
- local matrices include:
the local stiffness matrix A, which measures the alignment of the gradients of two basis functions  i.e. how much the element resist deformation (spring force) 
the local mass matrix M, which measures how much basis functions occupy the same region of space (overlap) i.e. how much the triangle resist acceleration


\subsection{Global assembly}
global assembly went through a straightforward algorithm. Each vertex of the triangle has its own global coordinates, so the idea is to put the local matrix in the similar position in global as it has in the domain.


\subsection{Dirichlet boundary conditions}
- all the nodes on the boundary will vanish


\subsection{Eigenvalue solver}


\subsection{Eigenvalues}
for P1
- the mesh refinement was reducing the error by half 
- fist eigenvalue was approximated quite precisely

\subsection{Eigenvectors}

\subsection{Degenerate eigenspaces}
for 2D case second and third eigenvalues were the same, this can be explained by the exact eigenfunctions which is $u(x, y) = sin(2*\pi*x)*sin(\pi*y)$ and $u(x, y) = sin(\pi*x)*sin(2*\pi*y)$

\subsection{Error analysis}


\section{Convergence analysis}
the errors were converging quadratically which can be concluded by the graphs of convergence, the convergence rates of the eigenvalues and eigenvectors matched the expected theoretical convergence rates


\section{Numerical results for the Laplace Eigenproblem}
- mesh refinement
- convergence plots
- eigenfunctions
- exact vs fem
- degenerate eigenspaces


\section{
    \text{Numerical results for the Schr\"{o}dinger's Eigenproblem}
}
- constant potential
- potential matrix
- harmonic oscillator and coulomb
- comparison with Laplace



\section{Discussion}
- why stiffness matrix alone is insufficient
- how potential modifies the operator
- why constant potential shifts the eigenvalue
- why varying potentials also change eigenfucntions
- why validation becomes harder

\section{Conclusion}


\begin{thebibliography}{00}
\bibitem{b1} G. Eason, B. Noble, and I. N. Sneddon, ``On certain integrals of Lipschitz-Hankel type involving products of Bessel functions,'' Phil. Trans. Roy. Soc. London, vol. A247, pp. 529--551, April 1955.
\bibitem{b2} J. Clerk Maxwell, A Treatise on Electricity and Magnetism, 3rd ed., vol. 2. Oxford: Clarendon, 1892, pp.68--73.
\bibitem{b3} I. S. Jacobs and C. P. Bean, ``Fine particles, thin films and exchange anisotropy,'' in Magnetism, vol. III, G. T. Rado and H. Suhl, Eds. New York: Academic, 1963, pp. 271--350.
\bibitem{b4} K. Elissa, ``Title of paper if known,'' unpublished.
\bibitem{b5} R. Nicole, ``Title of paper with only first word capitalized,'' J. Name Stand. Abbrev., in press.
\bibitem{b6} Y. Yorozu, M. Hirano, K. Oka, and Y. Tagawa, ``Electron spectroscopy studies on magneto-optical media and plastic substrate interface,'' IEEE Transl. J. Magn. Japan, vol. 2, pp. 740--741, August 1987 [Digests 9th Annual Conf. Magnetics Japan, p. 301, 1982].
\bibitem{b7} M. Young, The Technical Writer's Handbook. Mill Valley, CA: University Science, 1989.
\end{thebibliography}
\vspace{12pt}
\end{document}